\subsection{مبانی نظری و فرمول‌بندی ریاضی}
فرض کنید تانسور حالت‌های پنهان ورودی $\mathbf{X} \in \mathbb{R}^{B \times N \times D}$ باشد. پروجکشن‌های پرس‌وجو، کلید و مقدار با استفاده از لایه‌های خطی سازگار با آداپتورهای کم‌مرتبه (\lr{LoRALinear}) به صورت زیر استخراج می‌شوند:
\begin{equation}
\mathbf{Q} \in \mathbb{R}^{B \times H \times N \times d}, \quad \mathbf{K}, \mathbf{V} \in \mathbb{R}^{B \times H_{\text{kv}} \times N \times d}
\end{equation}
در صورت عدم تقارن تعداد سرهای پرس‌وجو و کلید/مقدار (\lr{GQA})، تانسورهای $\mathbf{K}$ و $\mathbf{V}$ با فاکتور تکرار $R = H / H_{\text{kv}}$ به ابعاد متناظر با $\mathbf{Q}$ گسترش می‌یابند.

ماتریس مجاورت گرافی $\mathbf{A} \in \{0, 1\}^{N \times N}$ بر اساس فاصله مکانی توکن‌ها و گره‌های لنگر به صورت زیر تعریف می‌شود:
\begin{equation}
\mathbf{A}_{i, j} = \begin{cases} 
1 & \text{if } |i - j| \le 16 \lor i \in \{0, N-1\} \lor j \in \{0, N-1\} \\
0 & \text{otherwise}
\end{cases}
\end{equation}
در مدل‌های دیکودر خودبازگشتی، ماسک علّی بر روی گراف اعمال می‌شود ($\mathbf{A}_{i,j}^{\text{causal}} = \mathbf{A}_{i,j} \land (i \ge j)$). امتیازات توجه گرافی با مقیاس‌بندی $\frac{1}{\sqrt{d}}$ و ماسک‌گذاری مواضع نامعتبر به دست می‌آیند:
\begin{equation}
\mathbf{S}_{\text{graph}} = \operatorname{masked\_fill}\left(\frac{\mathbf{Q} \mathbf{K}^\top}{\sqrt{d}}, \neg \mathbf{A}, -\infty\right), \quad \mathbf{O}_{\text{graph}} = \operatorname{softmax}(\mathbf{S}_{\text{graph}}) \mathbf{V}
\end{equation}

در شاخه سراسری، با اعمال نگاشت غیرخطی نامنفی $\phi(x) = \operatorname{ELU}(x) + 1$، توجه خطی در حالت انکودر دوطرفه با تغییر ترتیب ضرب ماتریسی محاسبه می‌شود:
\begin{equation}
\mathbf{M}_{\text{kv}} = \phi(\mathbf{K})^\top \mathbf{V} \in \mathbb{R}^{B \times H \times d \times d}, \quad \mathbf{O}_{\text{global}} = \frac{\phi(\mathbf{Q}) \mathbf{M}_{\text{kv}}}{\phi(\mathbf{Q}) \sum_{j=1}^N \phi(\mathbf{K}_j)^\top + \epsilon}
\end{equation}
و در حالت علّی ديكودر از جمع تجمعی ضرب خارجی توکن‌ها در راستای زمان استفاده می‌گردد. در پایان، شبکه دروازه بردار وزن‌های تلفیق را تولید کرده و خروجی نهایی شکل می‌گیرد:
\begin{equation}
[\mathbf{g}_{\text{graph}}, \mathbf{g}_{\text{global}}] = \sigma(\mathbf{X} \mathbf{W}_{\text{gate}} + \mathbf{b}_{\text{gate}}), \quad \mathbf{O} = \mathbf{g}_{\text{graph}} \odot \mathbf{O}_{\text{graph}} + \mathbf{g}_{\text{global}} \odot \mathbf{O}_{\text{global}}
\end{equation}
